\chapter{Introducción}
\label{sec:intro}

Los Sistemas Ciberfísicos (CPS) integran estrechamente las capacidades de computación, comunicación y control para interactuar con procesos físicos en el mundo real. Asegurar la estabilidad de estos sistemas ante perturbaciones externas o incertidumbres en los parámetros es uno de los desafíos más críticos en la ingeniería. \rev{Las consecuencias de una estabilización deficiente están ampliamente documentadas: el apagón del noreste de Norteamérica de 2003, donde una combinación de fallas locales y lazos de control mal coordinados derivó en una inestabilidad de tensión en cascada que dejó sin suministro a unos 50 millones de personas \cite{uscanada2004blackout}; o los escenarios de control en red sometidos a ataques de inyección de datos falsos y secuestro de actuadores, capaces de desestabilizar la planta física pese a que cada componente parezca operar dentro de sus márgenes nominales \cite{cardenas2008secure, pasqualetti2013attack}.} En el ámbito de los sistemas dinámicos, el concepto de estabilidad —particularmente en el sentido de Lyapunov— se define fundamentalmente como la propiedad de un sistema de retornar a un estado de equilibrio o permanecer confinado dentro de una región segura tras ser sometido a una perturbación. Tradicionalmente, el diseño de controladores con garantías formales de estabilidad se aborda utilizando modelos matemáticos del sistema físico y postulando funciones candidatas de Lyapunov, las cuales actúan como certificados matemáticos de que la energía del sistema se disipa continuamente con el tiempo para volver a un estado de equilibrio.

Utilizando herramientas de la teoría de control clásico, como el Procedimiento S, estas ecuaciones dinámicas y sus derivadas temporales no lineales pueden ser reestructuradas. El propósito fundamental de esta reestructuración es simplificar condiciones de diseño algebraicamente complejas e infinitas restricciones acotadas por perturbaciones, transformándolas en un formato estandarizado y matemáticamente manejable. Dicha reestructuración transforma las restricciones originales de estabilidad y robustez en Desigualdades Matriciales Lineales (LMIs, por sus siglas en inglés). 

Si bien formular la síntesis del control como un problema de optimización bajo restricciones LMI es un enfoque matemático poderoso y con sólidas garantías teóricas, resolver estos problemas en cada instante de tiempo mediante métodos convencionales (como los solucionadores numéricos basados en algoritmos de punto interior) resulta computacionalmente muy costoso. Esto es especialmente limitante en el contexto de los CPS modernos, donde las condiciones de operación varían dinámicamente y las decisiones de control y optimización deben tomarse en tiempo real en cuestión de milisegundos.

Recientemente, los avances en aprendizaje automático y el diseño de leyes de control aprendidas han permitido abordar este problema. Por ejemplo, existen propuestas recientes orientadas al aprendizaje de extremo a extremo para generar un arranque en caliente (\textit{warm-start}) que acelera la optimización cuadrática en tiempo real \cite{sambharya2023end}, así como métodos basados en redes de proyección ortogonal y parametrizaciones convexas que aseguran un control seguro con certificados de estabilidad y seguridad aprendidos directamente de los datos \cite{grontas2026pinet, tordesillas2023rayen}.

En este contexto, la presente memoria de título aborda este problema mediante la arquitectura de redes neuronales \textit{LMI-Net}, propuesta originalmente por \rev{Tang, Goertzen y Azizan \cite{tang2026lminet}}. Esta arquitectura incorpora una capa de proyección diferenciable basada en la división de Douglas--Rachford, la cual fuerza de manera estricta y por construcción el cumplimiento de restricciones matriciales duras directamente sobre las salidas de la red. \rev{El aporte de esta memoria consiste en adaptar dicha arquitectura ---demostrada originalmente sobre sistemas lineales invariantes en el tiempo--- al problema de la estabilización robusta de sistemas LPV politópicos, validándola empíricamente sobre una base de datos estándar del área y caracterizando sus alcances y limitaciones.} El enfoque resultante permite optimizar criterios de desempeño dinámicos siguiendo una descomposición \textit{offline--online}: el grueso del cómputo se traslada a una fase de entrenamiento previa, de modo que la evaluación de un nuevo sistema se reduce a una inferencia de la red seguida de la proyección, sin violar bajo ninguna circunstancia las restricciones de estabilidad que exigen las LMIs.